\documentclass[11pt]{article}
\usepackage[margin=1in]{geometry}
\usepackage{parskip}
\usepackage{titlesec}
\usepackage{xcolor}
\usepackage{amsmath}

\definecolor{slidecolor}{RGB}{12,35,64}
\definecolor{caleb}{RGB}{0,100,60}

\titleformat{\section}{\large\bfseries\color{slidecolor}}{}{0em}{}[\vspace{-4pt}\rule{\linewidth}{0.6pt}]
\titlespacing{\section}{0pt}{16pt}{6pt}

\title{\textbf{Group Script --- Bias Collapse in Shallow ReLU Networks}\\
       \large NSF RTG Project 4 $\,\cdot\,$ Research Workshop $\,\cdot\,$ June 2026}
\author{NSF RTG Project 4}
\date{}

\begin{document}
\maketitle
\thispagestyle{empty}

{\small\textit{Note: Slides 6--15 cover Open Problem 4.1.}}

% ================================================================
\section*{Slide 1 \quad Title}

Welcome. This talk presents numerical results for three open problems on bias collapse in shallow ReLU networks.
The central phenomenon is that gradient flow causes the bias parameters of an overparameterized network to spontaneously organize into a small number of clusters, with no regularization required.
We will describe the theoretical setting, then walk through what we found numerically for each open problem.

% ================================================================
\section*{Slide 2 \quad Motivation}

The starting point is the overparameterization puzzle.
Modern neural networks have far more parameters than the problem they are solving would seem to require, and classical learning theory predicts this should cause overfitting.
In practice, it does not.
Gradient flow appears to self-simplify: it discards redundant parameters and finds low-complexity solutions.

Bias collapse is one concrete mechanism for this.
As an illustrative example: a network with $m = 1500$ neurons trained on $\sin(4\pi x)$ ends up with an effective complexity of $k = 7$ neurons --- the rest cluster together and act redundantly.
The number of surviving bias clusters equals the number of inflection points of the target function, which is a property of the target's structure, not the network's size.

This talk addresses three open problems that ask whether this phenomenon can be characterized precisely.
Does $C(m, f^*) \to k$ as $m \to \infty$?
Does the same collapse occur in higher dimensions?
And after collapse, is there a provable $L^2$ pruning bound?

% ================================================================
\section*{Slide 3 \quad Model and Gradient Flow}

The model is a 1D shallow ReLU network:
\[
  f(x) = \sum_{j=1}^m a_j\,\sigma(x - b_j), \qquad \sigma = \max(0,\cdot), \qquad x \in [-1,1].
\]
Each neuron places a kink at bias $b_j$ with slope change $a_j$.
We train by minimizing the $L^2$ loss $\mathcal{L} = \frac{1}{2}\int_{-1}^1 (f - f^*)^2\,dx$ under continuous gradient flow, which gives the ODEs
\[
  \dot{a}_j = -\int_{-1}^1 (f - f^*)\,\sigma(x - b_j)\,dx, \qquad
  \dot{b}_j = a_j \int_{b_j}^1 (f - f^*)\,dx.
\]
We solve these numerically via \texttt{scipy} RK45 up to time $T = 500$, with stationarity criterion $\max_j |\dot{a}_j|, |\dot{b}_j| < 0.01$.

The key quantity on the right is the integrated residual $R_j := \int_{b_j}^1(f - f^*)\,dx$, which appears in the bias ODE.
At a stationary point, every active neuron must satisfy $R_j = 0$.
This connects to the theory of Open Problem 4.1: a $k$-cluster arrangement centered at the inflection points of $f^*$ satisfies this condition exactly, because each inflection point is where the signed curvature of $f^*$ changes, which is where the residual integral changes sign.

% ================================================================
\section*{Slide 4 \quad The Three Open Problems}

We investigate three open problems.

OP 4.1 asks whether $C(m, f^*) \to k$ as $m \to \infty$, where $k$ is the number of sign-changing inflection points of $f^*$.
OP 4.2 asks whether an analogous collapse occurs for higher-dimensional targets, where the relevant structure is the angular distribution of neuron normals rather than bias positions on the real line.
OP 4.3 asks whether collapse provides a provable $L^2$ pruning guarantee: after the biases have clustered, can you replace each cluster with a single representative neuron and bound the resulting approximation error?

Each problem is introduced in detail immediately before its results.

% ================================================================
\section*{Slide 5 \quad Experimental Design}

For OP 4.1 we chose nine target functions: $\sin(n\pi x)$ for $n = 1$ through $7$, giving $k = 1, 3, 5, 7, 9, 11, 13$, plus two polynomial targets $x^3$ and $x^5 - 3x^3$ with $k = 1$ and $k = 3$ respectively.
The trigonometric family provides a clean closed form for the inflection locations and a wide range of $k$ values.
The polynomial targets are deliberate controls: they have the same $k$ as two trig targets but irregular inflection spacing, letting us isolate the effect of geometry.

The main sweep ran each target at $m \in \{50, 100, 250, 500, 1000, 2000, 3500, 5000\}$ at $T = 500$, for 72 completed runs, 70 of which are fully stationary.

Supporting experiments include a discrete gradient descent comparison, a lottery ticket survival analysis, pruning bound verification across all 72 runs, a 2D collapse experiment, and a 36-job instability test near $k$.

% ================================================================
\section*{Slide 6 \quad Open Problem 4.1: Cluster Count Conjecture}

Open Problem 4.1 is the cluster count conjecture.
The conjecture states that $\lim_{m\to\infty} C(m, f^*) = k$, where $k$ counts the sign-changing zeros of $(f^*)''$ in $(-1,1)$ and $C(m, f^*)$ is the number of distinct bias clusters at stationarity.

The reason we focus on bias clusters comes directly from the ODE structure.
The fixed-point condition $R_j = 0$ is satisfied exactly by a $k$-cluster arrangement at the inflection locations of $f^*$, so $k$ is not an arbitrary target --- it is the natural fixed-point count baked into the dynamics.

We chose continuous gradient flow rather than discrete gradient descent deliberately.
The ODE setting gives us exact stationarity conditions to verify and lets us distinguish genuine fixed points from slowly drifting trajectories.
As a sanity check we also ran discrete GD on the same targets, and $C$ does not converge to $k$ under discrete GD for any target except $\sin(\pi x)$, confirming the continuous-time structure is essential.

The difficulty that makes this an open problem is that both above-$k$ and below-$k$ stationary states exist.
The conjecture requires showing gradient flow selects $C = k$ from generic initializations --- not merely that a $k$-cluster fixed point exists, but that it has a large enough basin to be reached reliably.

% ================================================================
\section*{Slide 7 \quad $C(m,f^*)$ Across All 9 Targets}

This is the main experimental result: 72 simulations across nine targets and eight values of $m$, all at $T = 500$.

The choice of targets was deliberate.
The trig family gives a wide range of $k$ with known inflection locations for direct comparison.
The two polynomial targets test whether $k$ alone determines convergence, or whether the geometry of the inflection points matters independently.

Every target exhibits the same qualitative non-monotone behavior: $C$ rises above $k$, peaks, then falls.
For low-$k$ trigonometric targets the fall reaches $C = k$ within our $m$ range, confirming the conjecture in those cases.
Polynomial targets and high-$k$ targets lag considerably.
The fact that rise-then-fall is universal across all nine function classes suggests it reflects something structural about the ODE, not a coincidence of any particular target.

$\sin(7\pi x)$ with $k = 13$ reaches $C = 13$ for the first time at $m = 5000$, with a stationarity caveat to come.

% ================================================================
\section*{Slide 8 \quad The Rise-Then-Fall Dynamic}

The pattern follows from what the plot shows.
At small $m$, insufficient neurons to populate all $k$ attractors simultaneously --- $C$ overshoots.
As $m$ grows, redundancy enables collapse and $C$ falls back.
The peak shifts to larger $m$ as $k$ grows, and the next slides illustrate this concretely.

% ================================================================
\section*{Slide 9 \quad Confirmation: $\sin(\pi x)$, $k = 1$}

The cleanest case.
All 1000 bias trajectories collapse to $x = 0$.
$C = 1 = k$ holds for every $m$ from 1000 to 5000, fully stationary throughout.
This is the only target confirmed robustly across the full $m$ sweep.

% ================================================================
\section*{Slide 10 \quad Confirmation: $\sin(3\pi x)$, $k = 5$}

For $k = 5$, five clusters align on the five inflection points at $m = 1000$.
Path: $C = 8 \to 14 \to 5$.
The apparent collapse to $C = 1$ at larger $m$ is a dense-packing artifact, not a structural change.

% ================================================================
\section*{Slide 11 \quad New at $m=5000$: $\sin(7\pi x)$, $k = 13$}

First instance of $C = 13 = k$ for the highest-$k$ target tested.
Path: $14 \to 36 \to 13$.
Caveat: $\max|\dot{a}_j| = 0.034 > 0.01$, so stationarity is unconfirmed and the result may be transient.
We are re-running at $T = 1000$ to confirm.

% ================================================================
\section*{Slide 12 \quad The Below-$k$ Phenomenon}

In several runs we found fully stationary states with $C$ strictly below $k$.
This is the primary challenge to the conjecture: if gradient flow rests at a below-$k$ fixed point, the limit cannot be $k$ in general.

The critical question is whether these states are universal or initialization-dependent.
A single run cannot distinguish the two cases, so we selected the three clearest below-$k$ instances and re-ran each with two additional random seeds.

$C$ varies across seeds in every case.
$\sin(4\pi x)$ at $m = 2000$ gives $C \in \{2, 3, 5\}$ depending on the seed.
$x^5 - 3x^3$ at $m = 5000$ gives $\{2, 2, 3\}$ --- seed 25 reaches $C = k = 3$, confirming the $k$-cluster attractor exists for polynomial targets, just with a narrow basin.

Below-$k$ states are initialization-dependent local minima, not universal attractors.
This is a concrete constraint on any future proof: it cannot argue $C \to k$ unconditionally without accounting for the basin structure.

% ================================================================
\section*{Slide 13 \quad Polynomial vs.\ Trigonometric Targets}

The polynomial targets were included to test whether $k$ alone determines convergence, or whether the geometry of the inflection points matters independently.
The answer is clear.
$x^3$ and $\sin(\pi x)$ both have $k = 1$, yet $x^3$ reaches only $C = 4$ at $m = 5000$ while $\sin(\pi x)$ confirms $C = 1$ by $m = 1000$.

For $\sin(n\pi x)$ the inflection points are evenly spaced across $[-1,1]$, creating symmetric, balanced gradient attractors that neurons localize around early in training.
For $x^3$ there is one inflection point at the origin but the curvature is shallow --- far more redundancy is needed before collapse concentrates neurons there.

Any proof of OP 4.1 that treats inflection points as interchangeable --- depending only on $k$ and not on their spacing or the curvature profile --- will not account for the polynomial behavior.

% ================================================================
\section*{Slide 14 \quad Stationarity and Fixed-Point Verification}

A key methodological concern is whether runs that appear converged are genuine fixed points or simply moving slowly.
We addressed this by substituting the final $(a_j, b_j)$ back into the ODEs and checking the velocities directly.

For $\sin(\pi x)$ at $m = 1000$, amplitude and bias velocities are in the range $10^{-5}$ to $10^{-3}$ --- negligible.
The right panel shows $R_j \approx 0$ for every active neuron, which is not just a smallness check but the theoretical fixed-point condition from the ODE, verified to hold numerically.
The stationarity threshold $\max|\dot{a}_j| < 0.01$ used throughout the sweep therefore corresponds to the regime where $R_j \approx 0$ is already well-satisfied --- it is a principled criterion, not an arbitrary cutoff.

The discrete GD comparison reinforces why we used gradient flow.
Under discrete GD with a comparable training budget, $C$ does not converge to $k$ for any target except $\sin(\pi x)$.
The fixed-point structure that drives collapse is a feature of the continuous-time limit and breaks down under discretization.

% ================================================================
\section*{Slide 15 \quad Instability Near $k$: Results (36/36 complete)}

This experiment was designed to test the mechanism by which $C$ could approach $k$.
There are two natural ways gradient flow could drive $C$ toward $k$: above-$k$ states could be unstable and dissolve, or below-$k$ states could be unstable and grow.
We tested both, and the results rule out both mechanisms.

For the eight above-$k$ runs we continued the ODE for $T_{\text{perturb}} = 1000$ additional time units.
$C$ did not change in any case.
Above-$k$ states are stable ODE fixed points; they do not spontaneously dissolve toward $k$.

For below-$k$ runs we injected one extra neuron with amplitude $0.01$, placed either near an existing cluster or maximally distant from all clusters.
In all twelve cases the injected neuron's amplitude remained at $0.01$ throughout --- the ODE provides zero net driving force for a new cluster to form.

For exact-$k$ runs, the system returned to $C = k$ in all sixteen cases after perturbation, confirming exact-$k$ configurations are stable attractors.

The three-category design characterizes the full local stability picture around $k$.
The consequence for OP 4.1.1, which asks for a proof that $C(m, f^*) \leq k$ independently of $m$, is direct: the natural proof strategy would be to show above-$k$ states are unstable, but they are not.
A proof of the upper bound cannot come from the dynamics alone; it must come from initialization geometry --- showing that generic initializations cannot reach above-$k$ fixed points at large $m$.

% ================================================================
\section*{Slide 16 \quad Open Problem 4.2: Higher-Dimensional Collapse}

Open Problem 4.2 asks whether an analogous collapse occurs for higher-dimensional targets.
In $d = 2$, each neuron defines an oriented hyperplane parameterized by an angle $\theta_j = \angle(w_j)$ and a normalized offset $\beta_j = b_j / \|w_j\|$.
Collapse in this setting means the $(\theta_j, \beta_j)$ cloud concentrates onto a low-dimensional subset of the hyperplane parameter space.

We chose three targets to provoke structurally distinct collapse modes.
The ridge target $\sin(2\pi u^\top x) + 0.5\sin(4\pi u^\top x)$ with $u = (3/5, 4/5)$ varies only along one oblique direction --- all normals should concentrate toward $\theta \approx 53°$.
The separable target $\sin(2\pi x_1) + 0.3\sin(2\pi x_2)$ has two active directions with unequal weights, predicted to produce two direction families of unequal size.
The radial target $\cos(\pi\|x\|)$ is rotationally symmetric --- no preferred direction, so entropy should stay near uniform.

The main diagnostic is angular entropy $H(\{\theta_j\})$: it is near zero if all normals align, and near $\log(36) \approx 3.58$ if directions are uniformly spread.
We swept $m \in \{64, 256, 512, 1024\}$ over 6000 epochs with Adam.

% ================================================================
\section*{Slide 17 \quad OP 4.2: Results}

The summary grid shows entropy, cluster count, and prune ratio for all three targets across widths.

At small width, collapse mode tracks target structure.
Ridge shows the strongest directional concentration; separable shows two-family splitting; radial resists collapse as predicted.

The surprising finding is that collapse weakens monotonically as $m$ grows, opposite to what the large-$m$ theoretical prediction would suggest.
At $m = 1024$, radial reaches $H = 3.53$, essentially at the uniform maximum of 3.58, confirming the rotational symmetry prediction --- but only because the network is wide enough to cover all orientations without collapsing.
For ridge and separable, entropy increases with $m$ but collapse persists at small width.

The interpretation is that collapse in 2D is capacity-constrained.
A small network cannot fit the target without concentrating directions; a wide network can fit it while remaining spread.
This is an open question for future theoretical work.

% ================================================================
\section*{Slide 18 \quad Open Problem 4.3: Pruning Bound}

Open Problem 4.3 asks whether collapse provides a provable $L^2$ pruning guarantee.
After gradient flow, biases have formed $k$ clusters.
Pruning replaces each cluster with a single representative neuron, giving a network $\tilde{f}$ with only $k$ active neurons.
The bound from the slides is
\[
  \|\tilde{f} - f\|_{L^2} \leq \delta \cdot \sum_{j=1}^m |a_j|,
\]
where $\delta$ is the maximum intra-cluster diameter and $\sum|a_j|$ is the total weight mass.

The bound holds because within a cluster all biases $b_j$ differ by at most $\delta$, so $\sigma(x - b_j) \approx \sigma(x - b_{\text{center}})$ up to error $\delta$.
Summing over $m$ neurons, errors accumulate at most as $\delta \cdot \sum|a_j|$.

For the bound to be non-vacuous --- meaning it goes to zero as networks grow --- we need $\delta \cdot \sum|a_j| \to 0$.
Our data shows $\sum|a_j| \sim \sqrt{m}$, so the bound grows with $m$ unless $\delta$ shrinks fast enough.
$\delta$ does not shrink with $m$ at fixed $T$, but it does shrink with longer training time.
The path to a formal bound is therefore a study of $\delta$ as a function of $T$, not $m$.

% ================================================================
\section*{Slide 19 \quad OP 4.3: Pruning Bound Verification}

We verified the bound across all 72 runs.
It holds in every case, with tightness ratios --- actual error divided by bound --- ranging from 0.0002 to 0.148.
The bound holds with substantial slack throughout.

The key obstacle is visible in the plot: $\sum|a_j|$ grows roughly as $\sqrt{m}$, and $\delta$ does not shrink with $m$ at $T = 500$, so the bound loosens as $m$ grows even as the actual pruning error plateaus.
For $\sin(\pi x)$, the actual error stabilizes near 2.5 from $m = 1000$ onward, but the bound grows from roughly 35 to 88 over the same range.

Making the bound formally useful requires showing $\delta \to 0$ as $T \to \infty$ at fixed $m$, then combining that with the $\sqrt{m}$ growth rate to find a joint limit.

% ================================================================
\section*{Slide 20 \quad OP 4.1.4: Lottery Ticket --- Which Neurons Survive?}

This experiment addresses a supporting question for OP 4.1: which neurons survive gradient flow collapse to become cluster representatives, and why?

The setup exploits the fact that at $t = 0$ all amplitudes $a_j \sim \mathcal{N}(0, 0.01)$ are near zero, so the only structured quantity is the initial bias placement $b_j^{(0)}$.
A neuron whose initial bias happens to fall near an inflection point of $f^*$ is geometrically lucky.
We test whether such neurons preferentially survive.

We ran five training conditions per (target, $m$): a full $m$-neuron baseline; $k$ neurons selected by proximity to inflection points (\texttt{geometric}); the same $k$ bias positions but fresh random amplitudes (\texttt{bias\_only}); the same $k$ amplitudes but random bias positions (\texttt{amp\_only}); and $k$ randomly selected neurons (\texttt{random\_k}).

The results are clear.
\texttt{geometric} and \texttt{bias\_only} perform identically in every case --- initial amplitude values carry no information, only bias position matters.
\texttt{amp\_only} is consistently the worst condition, worse even than \texttt{random\_k}: good amplitudes paired with random positions perform worse than good positions paired with random amplitudes.

The performance claim --- that just the $k$ lucky neurons can match the full network --- is false: they achieve 100$\times$ to 50{,}000$\times$ higher loss.
$k$ neurons produce only $k$ kinks in the piecewise-linear approximation; the full network needs redundancy to achieve fine fits.
But the survival claim is supported: bias position at $t = 0$ is the sole informative quantity for which neurons become cluster representatives.

% ================================================================
\section*{Slide 21 \quad Summary of Evidence}

To summarize.
For OP 4.1, the conjecture $C \to k$ is confirmed for $\sin(\pi x)$ at all tested widths.
$\sin(7\pi x)$ reaches $C = 13 = k$ for the first time at $m = 5000$, stationarity pending.
Below-$k$ states are initialization-dependent local minima; the $k$-cluster basin exists but is not the only attractor.

The instability experiment is complete.
Above-$k$ states are stable fixed points; exact-$k$ states are stable attractors; below-$k$ injection-resistant.

For OP 4.2, collapse mode tracks target structure at small width, but weakens as width grows --- partial confirmation, with the large-$m$ behavior remaining an open question.

For OP 4.3, the pruning bound holds in all 72 runs.
Making it non-vacuous requires $\delta \to 0$ with $T$, which is not yet established.

For OP 4.1.4, bias position at $t = 0$ is the sole informative quantity for survival.

% ================================================================
\section*{Slide 22 \quad Open Questions and Next Steps}

The most immediate open question for OP 4.1 is whether $C = 13$ for $\sin(7\pi x)$ at $m = 5000$ is a genuine fixed point; we are running $T = 1000$ to check.
More broadly, a proof of $C \to k$ must account for initialization geometry --- below-$k$ local minima exist, so the argument must show the $k$-cluster basin dominates at large $m$, not that other basins are absent.
The result that above-$k$ states are stable rules out instability as the mechanism for the upper bound OP 4.1.1.

For OP 4.2, the open question is why collapse weakens at large $m$.
The theory predicts the opposite.
Collapse appears to be capacity-constrained, and a theory for this behavior does not yet exist.

For OP 4.3, the path to a formal proof runs through a study of $\delta(T)$ at fixed $m$: does $\delta \to 0$ as $T \to \infty$, and at what rate?
If $\delta = O(T^{-\gamma})$ for some $\gamma > 0$ and $\sum|a_j| = O(m^{1/2})$, then $\delta \cdot \sum|a_j| \to 0$ as $T \to \infty$ at fixed $m$, making the bound formally non-vacuous.

% ================================================================
\section*{Slide 23 \quad Conclusion}

To close.
Gradient flow is not only an optimizer; it appears to act as a compression algorithm that discovers the minimum-complexity representation of the target's intrinsic structure.
The inflection points of $f^*$ act as fixed-point attractors: neurons near one survive collapse, the rest dissolve.

We have provided numerical evidence for this picture across 72 runs on 9 targets, with three supporting experiments on stability, geometry, and pruning.
The conjecture is confirmed for low-$k$ targets and supported for the first time for $k = 13$.
The main obstacle to a complete proof is the initialization geometry question --- showing the $k$-cluster basin dominates generically.
A rigorous answer would give a mathematical explanation for why overparameterized networks generalize.

Thank you.

\end{document}
